% 定理环境3

\ifx\allfiles\undefined
\documentclass[12pt, a4paper, oneside, UTF8]{ctexbook}  % +  这一句是新增加的

% \documentclass[12pt, a4paper,oneside, UTF8]{ctexbook}
% \usepackage[dvipsnames]{xcolor}   % (环境2)删除
\usepackage[dvipsnames, svgnames]{xcolor} % (环境2)替换原来的 \usepackage[dvipsnames]{xcolor}
\usepackage[strict]{changepage} % (环境2)添加
\usepackage{framed}             % (环境2)添加
\usepackage{amsmath}   % 数学公式
\usepackage{amsthm}    % 定理环境
\usepackage{amssymb}   % 更多公式符号
\usepackage{graphicx}  % 插图
\usepackage{mathrsfs}  % 数学字体
\usepackage{enumitem}  % 列表
\usepackage{geometry}  % 页面调整
\usepackage{unicode-math}
\usepackage[colorlinks,linkcolor=black]{hyperref}

\usepackage{titling}    % 封面添加的

\usepackage{tcolorbox}  % (环境1)添加的
\tcbuselibrary{most}    % (环境1)添加的

\graphicspath{ {figure/},{../figure/}, {config/}, {../config/} }  % 配置图形文件检索目录
\linespread{1.5} % 行高

% 页码设置
\geometry{top=25.4mm,bottom=25.4mm,left=20mm,right=20mm,headheight=2.17cm,headsep=4mm,footskip=12mm}

% 设置列表环境的上下间距
\setenumerate[1]{itemsep=5pt,partopsep=0pt,parsep=\parskip,topsep=5pt}
\setitemize[1]{itemsep=5pt,partopsep=0pt,parsep=\parskip,topsep=5pt}
\setdescription{itemsep=5pt,partopsep=0pt,parsep=\parskip,topsep=5pt}



% % 原始定理环境
% % ########## 定理环境 start ##############################################################
% \theoremstyle{definition}
% \newtheorem{defn}{\indent 定义}[section]

% \newtheorem{lemma}{\indent 引理}[section]    % 引理 定理 推论 准则 共用一个编号计数
% \newtheorem{thm}[lemma]{\indent 定理}
% \newtheorem{corollary}[lemma]{\indent 推论}
% \newtheorem{criterion}[lemma]{\indent 准则}

\newtheorem{proposition}{\indent 命题}[section]
\newtheorem{example}{\indent \color{SeaGreen}{例}}[section] % 绿色文字的 例 ，不需要就去除\color{SeaGreen}{}
\newtheorem*{rmk}{\indent 注}

% % 两种方式定义中文的 证明 和 解 的环境：
% % 缺点：\qedhere 命令将会失效【技术有限，暂时无法解决】
\renewenvironment{proof}{\par\textbf{证明.}\;}{\qed\par}
\newenvironment{solution}{\par{\textbf{解.}}\;}{\qed\par}

% % 缺点：\bf 是过时命令，可以用 textb f等替代，但编译会有关于字体的警告，不过不影响使用【技术有限，暂时无法解决】
% %\renewcommand{\proofname}{\indent\bf 证明}
% %\newenvironment{solution}{\begin{proof}[\indent\bf 解]}{\end{proof}}
% % ######### 定理环境 end  ##############################################################





% % #### 将 config.tex 中的定理环境的对应部分替换为如下内容   (环境2)
% % 定义单独编号，其他四个共用一个编号计数 这里只列举了五种，其他可类似定义（未定义的使用原来的也可）
% \newtcbtheorem[number within=section]{defn}%
% {定义}{colback=OliveGreen!10,colframe=Green!70,fonttitle=\bfseries}{def}

% \newtcbtheorem[number within=section]{lemma}%
% {引理}{colback=Salmon!20,colframe=Salmon!90!Black,fonttitle=\bfseries}{lem}

% % 使用另一个计数器 use counter from=lemma
% \newtcbtheorem[use counter from=lemma, number within=section]{them}%
% {定理}{colback=SeaGreen!10!CornflowerBlue!10,colframe=RoyalPurple!55!Aquamarine!100!,fonttitle=\bfseries}{them}

% \newtcbtheorem[use counter from=lemma, number within=section]{criterion}%
% {准则}{colback=green!5,colframe=green!35!black,fonttitle=\bfseries}{cri}

% \newtcbtheorem[use counter from=lemma, number within=section]{corollary}%
% {推论}{colback=Emerald!10,colframe=cyan!40!black,fonttitle=\bfseries}{cor}
% % colback=red!5,colframe=red!75!black

% % 这个颜色我不喜欢
% %\newtcbtheorem[number within=section]{proposition}%
% %{命题}{colback=red!5,colframe=red!75!black,fonttitle=\bfseries}{cor}

% % .... 命题 例 注 证明 解 使用之前的就可以（全文都是这种框框就很丑了），也可以按照上述定义 ...






% #### 将 config.tex 中的定理环境的对应部分替换为如下内容   (环境3)
\definecolor{greenshade}{rgb}{0.90,1,0.92} 		% 绿色文本框，竖线颜色设为 Green
\definecolor{redshade}{rgb}{1.00,0.88,0.88}   		% 红色文本框，竖线颜色设为 LightCoral
\definecolor{brownshade}{rgb}{0.99,0.95,0.9} 		% 莫兰迪棕色，竖线颜色设为 BurlyWood
\definecolor{lilacshade}{rgb}{0.95,0.93,0.98}    	% 淡紫色，竖线颜色设为 Plum
\definecolor{orangeshade}{rgb}{1.00,0.88,0.82} 		% 橙色，竖线颜色设为 DarkOrange
\definecolor{lightblueshade}{rgb}{0.8,0.92,1} 		% 淡蓝色，竖线颜色设为 LightSkyBlue
\theoremstyle{definition}
\newtheorem{defn1}{\indent 定义}[section]
\newtheorem{lemma1}{\indent 引理}[section]
\newtheorem{thm1}[lemma1]{\indent 定理}
\newtheorem{corollary1}[lemma1]{\indent 推论}
\newtheorem{criterion1}[lemma1]{\indent 准则}
% \newtheorem*{rmk1}{\indent 注}

\newenvironment{formal}[2][]{%
    \def\FrameCommand{%
        \hspace{1pt}%
        {\color{#1}\vrule width 2pt}%
        {\color{#2}\vrule width 4pt}%
        \colorbox{#2}%
    }%
    \MakeFramed{\advance\hsize-\width\FrameRestore}%
    \noindent\hspace{-4.55pt}%
    \begin{adjustwidth}{}{7pt}\vspace{2pt}\vspace{2pt}}{%
        \vspace{2pt}\end{adjustwidth}\endMakeFramed%
}

% 定义
\newenvironment{defn}{\begin{formal}[Green]{greenshade}\begin{defn1}}{\end{defn1}\end{formal}}
% 定理
\newenvironment{thm}{\begin{formal}[LightSkyBlue]{lightblueshade}\begin{thm1}}{\end{thm1}\end{formal}}
% 引理
\newenvironment{lemma}{\begin{formal}[Plum]{lilacshade}\begin{lemma1}}{\end{lemma1}\end{formal}}
% 推论
\newenvironment{corollary}{\begin{formal}[BurlyWood]{brownshade}\begin{corollary1}}{\end{corollary1}\end{formal}}
% 准则
\newenvironment{criterion}{\begin{formal}[DarkOrange]{orangeshade}\begin{criterion1}}{\end{criterion1}\end{formal}}
% 注
% \newenvironment{rmk}{\begin{formal}[LightCoral]{redshade}\begin{rmk1}}{\end{rmk1}\end{formal}}






% ↓↓↓↓↓↓↓↓↓↓↓↓↓↓↓↓↓ 以下是自定义的命令  ↓↓↓↓↓↓↓↓↓↓↓↓↓↓↓↓

% 用于调整表格的高度  使用 \hline\xrowht{25pt}
\newcommand{\xrowht}[2][0]{\addstackgap[.5\dimexpr#2\relax]{\vphantom{#1}}}

% 表格环境内长内容换行
\newcommand{\tabincell}[2]{\begin{tabular}{@{}#1@{}}#2\end{tabular}}

% 使用\linespread{1.5} 之后 cases 环境的行高也会改变，重新定义一个 ca 环境可以自动控制 cases 环境行高
\newenvironment{ca}[1][1]{\linespread{#1} \selectfont \begin{cases}}{\end{cases}}
% 和上面一样
\newenvironment{vx}[1][1]{\linespread{#1} \selectfont \begin{vmatrix}}{\end{vmatrix}}

\def\d{\textup{d}} % 直立体 d 用于微分符号 dx
\def\R{\mathbb{R}} % 实数域
\newcommand{\bs}[1]{\boldsymbol{#1}}    % 加粗，常用于向量
\newcommand{\ora}[1]{\overrightarrow{#1}} % 向量

% 数学 平行 符号
\newcommand{\pll}{\kern 0.56em/\kern -0.8em /\kern 0.56em}

% 用于空行\myspace{1} 表示空一行 填 2 表示空两行  
\newcommand{\myspace}[1]{\par\vspace{#1\baselineskip}}

\begin{document}
% \renewcommand*{\maketitle}{
    \begin{titlepage}
        \newgeometry{margin = 0in}
        \parindent=0pt
        % 注意这里使用了一张图片，封面使用的图片放到 config文件夹下
        \includegraphics[width=\linewidth]{1.png}
        \vfill
        \begin{center}
            \parbox{0.618\textwidth}{
                \hfill {\bfseries \Huge \thetitle} \\[0.6pt]
                \rule{0.618\textwidth}{4pt}
            }
        \end{center}
        \vfill
        \begin{center}
            \parbox{0.618\textwidth}{
                \hfill\Large
                \begin{tabular}{r|}
                    作者：\theauthor \\
                    时间：\thedate   \\
                \end{tabular}
            }
        \end{center}
        \vfill
        \begin{center}
            \parbox[t]{0.7\textwidth}{\centering  }
        \end{center}
        \vfill
    \end{titlepage}
    \restoregeometry
    \thispagestyle{empty}
}

% 修改 名称 作者 时间
\title{REAL ANALYSIS}
\author{Tang}
\date{\today}
\thispagestyle{empty}
\maketitle

\thispagestyle{empty}
\begin{center}\Huge\textbf{前言}\end{center}
实分析笔记
\begin{flushright}
    \begin{tabular}{c}
        \today \\ 签名
    \end{tabular}
\end{flushright}

\newpage
\setcounter{page}{1}
\pagestyle{plain}
\pagenumbering{Roman}
\tableofcontents

\newpage
\pagestyle{plain}
\setcounter{page}{1}
\pagenumbering{arabic} % 单独编译时，其实不用编译封面目录之类的，如需要不注释这句即可
\else
\fi
%  ↓↓↓↓↓↓↓↓↓↓↓↓↓↓↓↓↓↓↓↓↓↓↓↓↓↓↓↓ 正文部分 ↓↓↓↓↓↓↓↓↓↓↓↓↓↓↓↓↓↓↓↓↓↓↓↓↓↓↓↓

\chapter{Lebesgue 测度}
\section{Jordan 测度}

\begin{defn}
    对每一$\R$中的紧集$E$, 定义
    \[ J^*(E)=inf\left \{\sum_{i}^{n}\ell(I_i):\left \{ I_i \right \}_{i\ge1}^{n}\mbox{是一列区间并且}E\subseteq\bigcup_{i=1}^{n}I_i\right \} \]
    称,$J^*(E)$为$E$的\textbf{\emph{Jordan 外测度}}.

    \[ J_*(E)=inf\left \{\sum_{i}^{n}\ell(I_i):\left \{ I_i \right \}_{i\ge1}^{n}\mbox{是一列区间并且}\bigcup_{i=1}^{n}I_i\subseteq E\right \} \]
    称,$J_*(E)$为$E$的\textbf{\emph{Jordan 内测度}}.

    若$J^*(E)=J_*(E)$, 称$E$为\textbf{\emph{Jordan 可测}}, $J(E)=J^*(E)=J_*(E)$为$E$的\textbf{\emph{Jordan 测度}}.
\end{defn}

\begin{rmk}
    Jordan 测度性质
    \begin{itemize}
        \item Jordan 测度定义为有限区间的逼近
        \item Jordan 可测集只构成代数
        \item Jordan 测度只满足有限可加性, 非可数可加
    \end{itemize}
\end{rmk}

\begin{example}
    设$C$是\textbf{\emph{Cantor}}三分集, 证$J(C)=0$
    \begin{proof}
        略.
    \par 易证$J^*(C) \geq J_*(C)$, 只需证$J^*(C)=0$即可.
    \end{proof}
\end{example}

\begin{example}
    $J_*(\mathbb{Q}\bigcap[0, 1])==0$ and $J^*(\mathbb{Q}\bigcap[0, 1])==1$
    \begin{proof}
        略.
    \end{proof}
\end{example}

\begin{defn}
    映射$\mu :\mathbb{R}\to [0, +\infty]$, $A$, $B$为$\mathbb{R}$中集合, 该映射满足
    \begin{itemize}
        \item $\mu(A)\leq \mu(B)$, 若$A\subseteq B$.
        \item $\mu(\bigcup_{i=1}^{\infty}A_i) =\sum_{i=1}^{\infty}\mu(A_i)$, 若$A_i$互不相交.
        \item $\mu([0,1])=1$.
        \item $\mu(E)=\mu(E+x)$, 对$\forall x\in\R$.
    \end{itemize}
    \par 若$\exists$ $\mu(E)$, 其中$E\subset \R$, 则称$E$为 \textbf{\emph{$\mu$-可测}} 集.
\end{defn}


%  ↑↑↑↑↑↑↑↑↑↑↑↑↑↑↑↑↑↑↑↑↑↑↑↑↑↑↑↑ 正文部分 ↑↑↑↑↑↑↑↑↑↑↑↑↑↑↑↑↑↑↑↑↑↑↑↑↑↑↑↑
\ifx\allfiles\undefined
\end{document}
\fi